\subsection{The original tau-base weight line in the exponential reflection flow}\label{endpoint-f1_reflection--the-original-tau-base-weight-line-in-the-exponential-reflection-flow}

13 September 2026. This is a calculation on the existing Split-Zero arithmetic construction, not a replacement finite-field model and not a weight estimate. It joins the original moment line in the tau-base right adjoint to an exact reflection and flow on the already constructed polynomial-exponential family. The finite transport identity extends through the special fibre even though the corresponding contour gauge contains an essential exponential in ENDPOINTSOURCE8AA488F2AE0FA16E20DC8118SLOT0END.

\subsubsection{1. Retained source and conventions}\label{endpoint-f1_reflection--retained-source-and-conventions}

The controlling source is \texttt{Tau\_Base\_Cohomology\_2026-09-12/NOTE.md}, SHA256 \texttt{d03e71ad18f187bd89880cb8ca6fc9c5d6f260b3de8e8e5702c27db97e7b63c3}. Its whole body was read. Sections 3--7 supply, respectively,

ENDPOINTSOURCE8AA488F2AE0FA16E20DC8118SLOT1END

The original action is ENDPOINTSOURCE8AA488F2AE0FA16E20DC8118SLOT2END. Its two source legs are ENDPOINTSOURCE8AA488F2AE0FA16E20DC8118SLOT3END, not two copies of ENDPOINTSOURCE8AA488F2AE0FA16E20DC8118SLOT4END. Its Mellin multiplier is ENDPOINTSOURCE8AA488F2AE0FA16E20DC8118SLOT5END. The actual second moment ENDPOINTSOURCE8AA488F2AE0FA16E20DC8118SLOT6END has multiplier ENDPOINTSOURCE8AA488F2AE0FA16E20DC8118SLOT7END; its k-fold tensor is the specified line ENDPOINTSOURCE8AA488F2AE0FA16E20DC8118SLOT8END with multiplier ENDPOINTSOURCE8AA488F2AE0FA16E20DC8118SLOT9END. The dual action is consequently ENDPOINTSOURCE8AA488F2AE0FA16E20DC8118SLOT10END times precomposition by ENDPOINTSOURCE8AA488F2AE0FA16E20DC8118SLOT11END. All these maps fix support masks and carry supported zero to supported zero, with external absence tau carried to tau.

Use an actual full quartet packet and its original cyclic tensor module,

ENDPOINTSOURCE8AA488F2AE0FA16E20DC8118SLOT12END

Here the retained quartet and its complete multiplicities give

ENDPOINTSOURCE8AA488F2AE0FA16E20DC8118SLOT13END

Indeed complex conjugation and reflection permute the tensor-sum roots with their original multiplicities; the leading coefficient of the reflected monic degree-q polynomial is ENDPOINTSOURCE8AA488F2AE0FA16E20DC8118SLOT14END. This statement is about those actual symmetry-stable packets. It does not invent an off-line zero or discard a nilpotent direction. The arithmetic inclusion remains ENDPOINTSOURCE8AA488F2AE0FA16E20DC8118SLOT15END. In particular its full Taylor unit is not absorbed into a scalar or removed from the comparison.

The existing exponential family uses the antiderivative with ENDPOINTSOURCE8AA488F2AE0FA16E20DC8118SLOT16END,

ENDPOINTSOURCE8AA488F2AE0FA16E20DC8118SLOT17END

Leading-term division makes H free of rank q over ENDPOINTSOURCE8AA488F2AE0FA16E20DC8118SLOT18END, with basis ENDPOINTSOURCE8AA488F2AE0FA16E20DC8118SLOT19END. In that basis the already proved connection is

ENDPOINTSOURCE8AA488F2AE0FA16E20DC8118SLOT20END

The matrix ENDPOINTSOURCE8AA488F2AE0FA16E20DC8118SLOT21END is multiplication by S in ENDPOINTSOURCE8AA488F2AE0FA16E20DC8118SLOT22END on these representatives. Multiplication by an arbitrary S-polynomial is not thereby asserted to descend to the full u-deformed differential cokernel. The connection and the chain maps below supply the actual descended operators.

\subsubsection{2. Reflection is a cochain map of the actual family}\label{endpoint-f1_reflection--reflection-is-a-cochain-map-of-the-actual-family}

Let ENDPOINTSOURCE8AA488F2AE0FA16E20DC8118SLOT23END and

ENDPOINTSOURCE8AA488F2AE0FA16E20DC8118SLOT24END

The matrix R on degree-below-q representatives has the explicit entries

ENDPOINTSOURCE8AA488F2AE0FA16E20DC8118SLOT25END

For every polynomial P, differentiation and (RW3) give

ENDPOINTSOURCE8AA488F2AE0FA16E20DC8118SLOT26END

In detail, the derivative term on the left is ENDPOINTSOURCE8AA488F2AE0FA16E20DC8118SLOT27END; its remaining term is ENDPOINTSOURCE8AA488F2AE0FA16E20DC8118SLOT28END, which is exactly the reflected right side. Thus degree-zero map ENDPOINTSOURCE8AA488F2AE0FA16E20DC8118SLOT29END and degree-one map R give a cochain isomorphism, with its inverse given by the same reflection. This extends over ENDPOINTSOURCE8AA488F2AE0FA16E20DC8118SLOT30END, where it is the original polynomial reflection on E. Adding coefficient conjugation gives the actual dagger operation, with parameter map ENDPOINTSOURCE8AA488F2AE0FA16E20DC8118SLOT31END.

On the companion matrices the same calculation is

ENDPOINTSOURCE8AA488F2AE0FA16E20DC8118SLOT32END

This follows either by reducing ENDPOINTSOURCE8AA488F2AE0FA16E20DC8118SLOT33END modulo ENDPOINTSOURCE8AA488F2AE0FA16E20DC8118SLOT34END, or by using the algebra isomorphism defined by (RW3). All entries and all repeated-root blocks remain; no diagonalization is used.

\subsubsection{3. The weight twist is the original potential's constant, not a choice}\label{endpoint-f1_reflection--the-weight-twist-is-the-original-potentials-constant-not-a-choice}

Integrating (RW3) with the original endpoint gives

ENDPOINTSOURCE8AA488F2AE0FA16E20DC8118SLOT35END ENDPOINTSOURCE8AA488F2AE0FA16E20DC8118SLOT36END

Neither ENDPOINTSOURCE8AA488F2AE0FA16E20DC8118SLOT37END nor ENDPOINTSOURCE8AA488F2AE0FA16E20DC8118SLOT38END is set to zero. Pull the target connection back along r, keeping u fixed while differentiating t. Since ENDPOINTSOURCE8AA488F2AE0FA16E20DC8118SLOT39END,

ENDPOINTSOURCE8AA488F2AE0FA16E20DC8118SLOT40END

Consequently the exact horizontal map on ENDPOINTSOURCE8AA488F2AE0FA16E20DC8118SLOT41END is

ENDPOINTSOURCE8AA488F2AE0FA16E20DC8118SLOT42END

Its logarithmic t-derivative is ENDPOINTSOURCE8AA488F2AE0FA16E20DC8118SLOT43END, so substituting (RW11) proves ENDPOINTSOURCE8AA488F2AE0FA16E20DC8118SLOT44END. This uses the full constant supplied by the original potential. It does not select a new norm on H.

For an oriented rapid-decay contour Gamma, its reflected contour is the pushforward ENDPOINTSOURCE8AA488F2AE0FA16E20DC8118SLOT45END with that orientation. The substitution ENDPOINTSOURCE8AA488F2AE0FA16E20DC8118SLOT46END has ENDPOINTSOURCE8AA488F2AE0FA16E20DC8118SLOT47END. Therefore the actual integrals satisfy

ENDPOINTSOURCE8AA488F2AE0FA16E20DC8118SLOT48END

The minus sign is the one-form pullback. The constant exponential is retained. The exponential identity also carries the original decay along Gamma to Gamma', so this substitution does not postulate a different contour class.

\subsubsection{4. A regular finite flow identity through the special fibre}\label{endpoint-f1_reflection--a-regular-finite-flow-identity-through-the-special-fibre}

The scaling calculation computes the regular lift of forward theta dilation from ENDPOINTSOURCE8AA488F2AE0FA16E20DC8118SLOT49END. Its coefficient transport C\_a is characterized by

ENDPOINTSOURCE8AA488F2AE0FA16E20DC8118SLOT50END

This linear matrix equation, with polynomial coefficients, has a unique entire solution in a, depending holomorphically on u,t. Its local power series recursion is explicit. It obeys

ENDPOINTSOURCE8AA488F2AE0FA16E20DC8118SLOT51END

To prove the product law, differentiate its right side in b and compare its initial value at b=0 with ENDPOINTSOURCE8AA488F2AE0FA16E20DC8118SLOT52END in (RW14). Uniqueness gives equality. In particular this finite transport is invertible even at u=0; there ENDPOINTSOURCE8AA488F2AE0FA16E20DC8118SLOT53END.

Its reflection is the new exact formula

ENDPOINTSOURCE8AA488F2AE0FA16E20DC8118SLOT54END

Proof: write the left side as B\_a. By (RW9),

ENDPOINTSOURCE8AA488F2AE0FA16E20DC8118SLOT55END

For ENDPOINTSOURCE8AA488F2AE0FA16E20DC8118SLOT56END, ENDPOINTSOURCE8AA488F2AE0FA16E20DC8118SLOT57END, the derivative of ENDPOINTSOURCE8AA488F2AE0FA16E20DC8118SLOT58END is ENDPOINTSOURCE8AA488F2AE0FA16E20DC8118SLOT59END. Thus the right side of (RW16) obeys this same equation and has the same initial value I. This proves (RW16) for all parameters, without dividing by u. It is not an attempt to evaluate the singular exponential (RW12) at u=0.

Along ENDPOINTSOURCE8AA488F2AE0FA16E20DC8118SLOT60END, the ratio of the two full scalar factors in (RW12) is exactly ENDPOINTSOURCE8AA488F2AE0FA16E20DC8118SLOT61END. At ENDPOINTSOURCE8AA488F2AE0FA16E20DC8118SLOT62END, forward theta dilation has moment line ENDPOINTSOURCE8AA488F2AE0FA16E20DC8118SLOT63END; inverse dilation U\_p has moment line ENDPOINTSOURCE8AA488F2AE0FA16E20DC8118SLOT64END. Thus the very same factor appears both in the actual tau-base adjunction and in the exact exponential-family reflection flow. It is not a finite-field cardinality assigned to the absolute base point.

Specializing (RW16) gives the full arithmetic identity

ENDPOINTSOURCE8AA488F2AE0FA16E20DC8118SLOT65END

Every block is retained:

ENDPOINTSOURCE8AA488F2AE0FA16E20DC8118SLOT66END

\subsubsection{5. The computed adjoint transport and its exact evaluation}\label{endpoint-f1_reflection--the-computed-adjoint-transport-and-its-exact-evaluation}

On coefficient sections the scaling operator is semilinear: ENDPOINTSOURCE8AA488F2AE0FA16E20DC8118SLOT67END. The k-fold moment line has ENDPOINTSOURCE8AA488F2AE0FA16E20DC8118SLOT68END. Therefore the actual twisted dual is

ENDPOINTSOURCE8AA488F2AE0FA16E20DC8118SLOT69END

Substitution proves, with no omitted scalar,

ENDPOINTSOURCE8AA488F2AE0FA16E20DC8118SLOT70END

The product rule (RW15) proves the composition law for (RW19). Its infinitesimal generator is

ENDPOINTSOURCE8AA488F2AE0FA16E20DC8118SLOT71END

At the actual special fibre and inverse dilation it becomes ENDPOINTSOURCE8AA488F2AE0FA16E20DC8118SLOT72END. This is precisely the action obtained by applying the right adjoint ENDPOINTSOURCE8AA488F2AE0FA16E20DC8118SLOT73END in the original source, then precomposing finite jet evaluation. In the conjugate packet, the source's dagger and residue pairing give its specified linear map to this twisted dual. No Hermitian positivity has been inferred from this bilinear calculation.

For completeness, the existing constant residue matrix

ENDPOINTSOURCE8AA488F2AE0FA16E20DC8118SLOT74END

has determinant ENDPOINTSOURCE8AA488F2AE0FA16E20DC8118SLOT75END and satisfies ENDPOINTSOURCE8AA488F2AE0FA16E20DC8118SLOT76END. Its t-independence follows because the difference of the two fractions starts at degree at most ENDPOINTSOURCE8AA488F2AE0FA16E20DC8118SLOT77END. With ENDPOINTSOURCE8AA488F2AE0FA16E20DC8118SLOT78END, the computed two-parameter pairing identity is

ENDPOINTSOURCE8AA488F2AE0FA16E20DC8118SLOT79END

Indeed transpose ENDPOINTSOURCE8AA488F2AE0FA16E20DC8118SLOT80END and apply the residue self-adjointness of ENDPOINTSOURCE8AA488F2AE0FA16E20DC8118SLOT81END. At t=0 this recovers the additive weight relation. For the actual arithmetic residue with denominator g, one must additionally retain the existing full unit ENDPOINTSOURCE8AA488F2AE0FA16E20DC8118SLOT82END and the actual arithmetic inclusion eta; the polynomial chi residue is not silently substituted for that pairing. No extension of that arithmetic unit as a globally invertible operator on every deformed fibre is claimed.

\subsubsection{6. Position in Deligne's argument and in the current program}\label{endpoint-f1_reflection--position-in-delignes-argument-and-in-the-current-program}

The printed primary text of Weil II, pp.~178, 202, 203 and 206, was read in the supplied S20 source-language records together with the existing primary-scan correction. In (3.3.5), the dual degree and Tate twist give ENDPOINTSOURCE8AA488F2AE0FA16E20DC8118SLOT83END. In (3.3.6), opposite bounds then constrain the image of the actual compact-support comparison. Page 203 (3.2.13) tensors the same eigenvalue with itself; its square is not replaced by an unspecified tensor constituent.

Our computation supplies a concrete part of the analogous tau-base diagram: the original moment-weight dual action has been transported through the exact exponential deformation and its reflection, including the full special-fibre blocks. It has not supplied Deligne's geometric upper weight bound for the theta quotient. The canonical arithmetic metrics G\_N and their original relation volumes remain unchanged.

The next use of this result is to transport the original residue-dual map and its arithmetic norm comparison along this explicit flow; the factor and parameter shift are now calculated, so they need not be guessed or replaced by a generic purity assertion. The external-product construction and signed determinant calculation are the continuations on those two sides.

\subsubsection{Provenance and verification scope}\label{endpoint-f1_reflection--provenance-and-verification-scope}

\begin{itemize}
\tightlist
\item
  Full controlling tau-base note as pinned above; no whole-paper Lean claim.
\item
  Existing exponential delivery \texttt{Tau\_Deligne\_Exponential\_Comparison/RESEARCH\_NOTE.md}, section 7 (37)--(42), read in full for the constant residue duality; DS66--69 supply the previously constructed family and original theta map.
\item
  The scaling calculation supplies the independent chain-level lift underlying (RW14). Formulas (RW6)--(RW13), (RW16), and their joining to the actual weight line are derived here, with their complete algebra above.
\item
  Exact finite fixtures are supporting checks, not replacements for these general proofs. No new Lean execution, RH proof, or remote publication is claimed by this note.
\end{itemize}
